%% ----------------------------------------------------------------------------
%%  §7.4  Derivate successive. Teorema di Schwarz
%% ----------------------------------------------------------------------------
\section{Derivate successive. Teorema di Schwarz}
\label{sec:schwarz}

Sia $f$ una funzione di due variabili definita in un intorno circolare $I_{\delta}$ di un
punto di $\R^{2}$, e supponiamo che in tutti i punti di $I_{\delta}$ esistano le derivate
parziali $f_{x}(x,y)$ e $f_{y}(x,y)$. Queste sono a loro volta funzioni di due
variabili, e possono ammettere derivate parziali:
\begin{equation}\label{eq:derivate-seconde-def}
  \frac{\partial}{\partial x}f_{x}, \qquad \frac{\partial}{\partial y}f_{x}, \qquad
  \frac{\partial}{\partial x}f_{y}, \qquad \frac{\partial}{\partial y}f_{y} .
\end{equation}
Esse si chiamano \emph{derivate parziali seconde} di $f$ e si indicano rispettivamente con
i simboli
\begin{equation}\label{eq:derivate-seconde-notazioni}
\begin{gathered}
  f_{xx}, \qquad f_{xy}, \qquad f_{yx}, \qquad f_{yy},\\
  \text{oppure}\qquad
  \frac{\partial^{2}f}{\partial x^{2}}, \quad \frac{\partial^{2}f}{\partial x\,\partial y}, \quad
  \frac{\partial^{2}f}{\partial y\,\partial x}, \quad \frac{\partial^{2}f}{\partial y^{2}} .
\end{gathered}
\end{equation}
Le $f_{xy}$, $f_{yx}$ si dicono \emph{derivate seconde miste}, le $f_{xx}$, $f_{yy}$
\emph{derivate seconde pure}. Nella notazione con i pedici l'ordine di derivazione si
legge da sinistra a destra: $f_{xy}=(f_{x})_{y}$ si ottiene derivando prima rispetto a $x$
e poi rispetto a $y$.

\begin{attenzione}[Una questione di convenzioni]
Nella notazione $\frac{\partial^{2}f}{\partial x\,\partial y}$ l'ordine di derivazione non è
uniforme nella letteratura: molti testi, interpretando il simbolo come
$\frac{\partial}{\partial x}\bigl(\frac{\partial f}{\partial y}\bigr)$, lo usano per $f_{yx}$.
Nel seguito useremo soprattutto i pedici, il cui significato è univoco; vedremo comunque
che, sotto ipotesi molto generali, le due derivate miste coincidono e la questione
diventa irrilevante.
\end{attenzione}

\begin{esempio}[Le derivate seconde di $x^{y}$]\label{es:derivate-xy}
La funzione $f(x,y)=x^{y}$, definita nel semipiano $D=\set{(x,y)\in\R^{2} : x>0}$, per $y$
fissato è una potenza, mentre per $x$ fissato è un esponenziale, che ai fini della
derivazione rispetto a $y$ conviene scrivere nella forma $f(x,y)=e^{y\log x}$. Le derivate
parziali prime valgono
\begin{equation}\label{eq:xy-prime}
  f_{x}=y\,x^{y-1}, \qquad f_{y}=x^{y}\log x ;
\end{equation}
le derivate seconde pure sono $f_{xx}=y(y-1)x^{y-2}$ e $f_{yy}=x^{y}(\log x)^{2}$. Per
ottenere $f_{xy}$ deriviamo $f_{x}$ rispetto a $y$ con la regola del prodotto, e per
$f_{yx}$ deriviamo $f_{y}$ rispetto a $x$:
\begin{equation}\label{eq:xy-miste}
\begin{aligned}
  f_{xy}&=x^{y-1}+y\,x^{y-1}\log x=x^{y-1}(1+y\log x),\\
  f_{yx}&=y\,x^{y-1}\log x+x^{y}\cdot\frac1x=x^{y-1}(y\log x+1).
\end{aligned}
\end{equation}
\end{esempio}

Il lettore avrà notato che, con passaggi intermedi differenti, si è ottenuto lo stesso
risultato. Ciò non è casuale ed è formalizzato nell'importante teorema che segue.

\begin{teorema}[di Schwarz]\label{teo:schwarz}\index{teorema!di Schwarz}
Sia $f$ una funzione definita in un intorno circolare $I_{\delta}$ di $(x_{0},y_{0})$, nel
quale esistono le derivate $f_{x}$, $f_{y}$, $f_{xy}$, $f_{yx}$. Se le derivate miste
$f_{xy}$ e $f_{yx}$ sono continue in $(x_{0},y_{0})$, allora
\begin{equation}\label{eq:schwarz}
  f_{xy}(x_{0},y_{0})=f_{yx}(x_{0},y_{0}) .
\end{equation}
\end{teorema}

\begin{schema}
\Ipotesi{\begin{itemize}[nosep, leftmargin=1.2em]
  \item $f_{x}$, $f_{y}$, $f_{xy}$, $f_{yx}$ esistono in $I_{\delta}(x_{0},y_{0})$;
  \item $f_{xy}$ e $f_{yx}$ sono continue in $(x_{0},y_{0})$.
\end{itemize}}
\Tesi{$f_{xy}(x_{0},y_{0})=f_{yx}(x_{0},y_{0})$.}
\end{schema}

\begin{strategia}
Consideriamo il rettangolo di vertici opposti $(x_{0},y_{0})$ e $(x,y)$
(\cref{fig:schwarz}) e la combinazione dei valori di $f$ nei quattro vertici
\[
  \Delta=f(x,y)-f(x,y_{0})-f(x_{0},y)+f(x_{0},y_{0}),
\]
una specie di «rapporto incrementale doppio». La stessa quantità $\Delta$ si può leggere
in due modi: come incremento, tra $x_{0}$ e $x$, della funzione $F(s)=f(s,y)-f(s,y_{0})$,
oppure come incremento, tra $y_{0}$ e $y$, della funzione $G(t)=f(x,t)-f(x_{0},t)$. Applicando
due volte il teorema di Lagrange nel primo modo si trova
$\Delta=f_{xy}(x_{1},y_{1})(x-x_{0})(y-y_{0})$; nel secondo modo
$\Delta=f_{yx}(x_{2},y_{2})(x-x_{0})(y-y_{0})$. Quindi il valore di $f_{xy}$ in un punto del
rettangolo coincide con quello di $f_{yx}$ in un altro punto del rettangolo; facendo tendere il rettangolo al punto
$(x_{0},y_{0})$, la continuità delle derivate miste fa il resto.
\end{strategia}

\begin{dimostrazione}
Sia $(x,y)\in I_{\delta}$ con $x\neq x_{0}$ e $y\neq y_{0}$. Il rettangolo chiuso di vertici
opposti $(x_{0},y_{0})$ e $(x,y)$ è contenuto in $I_{\delta}$: infatti ogni suo punto
$(s,t)$ verifica $\abs{s-x_{0}}\le\abs{x-x_{0}}$ e $\abs{t-y_{0}}\le\abs{y-y_{0}}$, e quindi
dista da $(x_{0},y_{0})$ non più di $(x,y)$. Definiamo
\begin{equation}\label{eq:schwarz-FG}
\begin{aligned}
  F(s)&=f(s,y)-f(s,y_{0}) \qquad (y \text{ fissato}),\\
  G(t)&=f(x,t)-f(x_{0},t) \qquad (x \text{ fissato}).
\end{aligned}
\end{equation}
Risulta
\[
\begin{aligned}
  F(x)-F(x_{0})&=f(x,y)-f(x,y_{0})-\bigl[f(x_{0},y)-f(x_{0},y_{0})\bigr],\\
  G(y)-G(y_{0})&=f(x,y)-f(x_{0},y)-\bigl[f(x,y_{0})-f(x_{0},y_{0})\bigr],
\end{aligned}
\]
per cui
\begin{equation}\label{eq:schwarz-uguaglianza}
  F(x)-F(x_{0})=G(y)-G(y_{0}) .
\end{equation}
La funzione $F$ è derivabile nell'intervallo di estremi $x_{0}$, $x$, con
$F'(s)=f_{x}(s,y)-f_{x}(s,y_{0})$. Per il teorema di Lagrange (\cref{teo:lagrange}) esiste
$x_{1}$ nell'intervallo di estremi $x_{0}$, $x$ tale che
\[
  F(x)-F(x_{0})=F'(x_{1})(x-x_{0})=\bigl[f_{x}(x_{1},y)-f_{x}(x_{1},y_{0})\bigr](x-x_{0}).
\]
Applicando di nuovo il teorema di Lagrange alla funzione $t\mapsto f_{x}(x_{1},t)$, che
è derivabile con derivata $f_{xy}(x_{1},t)$, nell'intervallo di estremi $y_{0}$, $y$,
otteniamo un punto $y_{1}$ di tale intervallo per cui
\begin{equation}\label{eq:schwarz-F}
  F(x)-F(x_{0})=f_{xy}(x_{1},y_{1})(x-x_{0})(y-y_{0}) .
\end{equation}
Procedendo in modo analogo con la funzione $G$ nell'intervallo di estremi $y_{0}$, $y$,
troviamo $y_{2}$ e poi $x_{2}$ tali che
\begin{equation}\label{eq:schwarz-G}
\begin{split}
  G(y)-G(y_{0})&=\bigl[f_{y}(x,y_{2})-f_{y}(x_{0},y_{2})\bigr](y-y_{0})\\
  &=f_{yx}(x_{2},y_{2})(x-x_{0})(y-y_{0}).
\end{split}
\end{equation}
Confrontando le \eqref{eq:schwarz-uguaglianza}, \eqref{eq:schwarz-F}, \eqref{eq:schwarz-G}
e dividendo per $(x-x_{0})(y-y_{0})\neq0$ si ottiene
\begin{equation}\label{eq:schwarz-punti}
  f_{xy}(x_{1},y_{1})=f_{yx}(x_{2},y_{2}),
\end{equation}
dove i punti $(x_{1},y_{1})$ e $(x_{2},y_{2})$ appartengono al rettangolo. Scegliamo ora
$(x,y)=(x_{0}+\frac rn,\,y_{0}+\frac rn)$, con $r>0$ fissato abbastanza piccolo e
$n\in\N$, $n\ge1$. I punti corrispondenti $(x_{1},y_{1})$, $(x_{2},y_{2})$ distano da
$(x_{0},y_{0})$ meno di $\frac{\sqrt2\,r}{n}$, e quindi tendono a $(x_{0},y_{0})$ per
$n\to+\infty$. Per la continuità di $f_{xy}$ e $f_{yx}$ in $(x_{0},y_{0})$
(\cref{oss:limite-piano-successioni}), passando al limite nella \eqref{eq:schwarz-punti}
si ottiene la tesi.
\end{dimostrazione}

\begin{figure}[htb]
\centering
\begin{tikzpicture}[scale=1.0]
  \draw[-Latex, grigio] (-0.3,0) -- (6.2,0) node[right, black] {$x$};
  \draw[-Latex, grigio] (0,-0.3) -- (0,3.6) node[above, black] {$y$};
  \draw[blu, dashed] (1.5,0.8) rectangle (5,3);
  \foreach \p/\lab/\pos in {(1.5,0.8)/{(x_{0},y_{0})}/below left,
      (5,0.8)/{(x,y_{0})}/below right, (1.5,3)/{(x_{0},y)}/above left,
      (5,3)/{(x,y)}/above right}
    {\fill[blu] \p circle (1.5pt); \node[font=\scriptsize, \pos] at \p {$\lab$};}
  \fill[rame] (3.9,1.6) circle (1.5pt) node[right, font=\scriptsize] {$(x_{1},y_{1})$};
  \fill[verde] (2.4,2.3) circle (1.5pt) node[right, font=\scriptsize] {$(x_{2},y_{2})$};
  \draw[tratteggio] (1.5,0) node[below, black, font=\scriptsize] {$x_{0}$} -- (1.5,0.8);
  \draw[tratteggio] (5,0) node[below, black, font=\scriptsize] {$x$} -- (5,0.8);
  \draw[tratteggio] (0,0.8) node[left, black, font=\scriptsize] {$y_{0}$} -- (1.5,0.8);
  \draw[tratteggio] (0,3) node[left, black, font=\scriptsize] {$y$} -- (1.5,3);
\end{tikzpicture}
\caption{I punti usati nella dimostrazione del teorema di Schwarz: i quattro vertici del
rettangolo, nei quali si calcola $f$, e i due punti interni $(x_{1},y_{1})$ e
$(x_{2},y_{2})$ forniti dal teorema di Lagrange, nei quali le derivate miste coincidono.}
\label{fig:schwarz}
\end{figure}

\begin{osservazione}[Le ipotesi di continuità non si possono omettere]
\label{oss:schwarz-controesempio}
Se le derivate seconde miste non sono continue, non necessariamente sono uguali.
Consideriamo la funzione
\begin{equation}\label{eq:schwarz-controesempio}
  g(x,y)=\frac{x^{3}y}{x^{2}+y^{2}} \quad\text{per }(x,y)\neq(0,0), \qquad g(0,0)=0 .
\end{equation}
La funzione è nulla sugli assi, quindi $g_{x}(0,0)=g_{y}(0,0)=0$. Per $(x,y)\neq(0,0)$,
\[
  g_{x}=\frac{3x^{2}y(x^{2}+y^{2})-2x^{4}y}{(x^{2}+y^{2})^{2}}, \qquad
  g_{y}=\frac{x^{3}(x^{2}+y^{2})-2x^{3}y^{2}}{(x^{2}+y^{2})^{2}} .
\]
In particolare $g_{x}(0,y)=0$ per ogni $y$, e $g_{y}(x,0)=x$ per ogni $x$. Quindi
\begin{align*}
  g_{xy}(0,0)&=\lim_{k\to0}\frac{g_{x}(0,k)-g_{x}(0,0)}{k}=0,\\
  g_{yx}(0,0)&=\lim_{h\to0}\frac{g_{y}(h,0)-g_{y}(0,0)}{h}=1 :
\end{align*}
le due derivate miste nell'origine sono diverse. Il teorema di Schwarz non è
contraddetto: le derivate miste di $g$ esistono in ogni punto, ma non sono continue
nell'origine.
\end{osservazione}

\begin{python}[Il teorema di Schwarz e il suo controesempio]
import sympy as sp

x, y, h = sp.symbols("x y h", real=True)
f = x**y                                 # definita per x > 0
fxy = sp.diff(f, x, y)              # prima rispetto a x, poi a y
fyx = sp.diff(f, y, x)              # prima rispetto a y, poi a x
print("f_xy =", sp.simplify(fxy), "   f_yx =", sp.simplify(fyx))
print("f_xy - f_yx =", sp.simplify(fxy - fyx))

# il controesempio: g = x^3 y / (x^2 + y^2), g(0, 0) = 0
g = x**3 * y / (x**2 + y**2)
gx0y = sp.simplify(sp.diff(g, x).subs(x, 0))   # g_x(0, y), y != 0
gyx0 = sp.simplify(sp.diff(g, y).subs(y, 0))   # g_y(x, 0), x != 0
print("g_x(0, y) =", gx0y, "   g_y(x, 0) =", gyx0)
# g_x(0,0) = g_y(0,0) = 0 (g è nulla sugli assi), quindi:
print("g_xy(0,0) =", sp.limit(gx0y.subs(y, h) / h, h, 0),
      "   g_yx(0,0) =", sp.limit(gyx0.subs(x, h) / h, h, 0))
\end{python}
\begin{risultato}
f_xy = x**(y - 1)*(y*log(x) + 1)    f_yx = x**(y - 1)*(y*log(x) + 1)
f_xy - f_yx = 0
g_x(0, y) = 0    g_y(x, 0) = x
g_xy(0,0) = 0    g_yx(0,0) = 1
\end{risultato}

Il teorema di Schwarz fornisce informazioni anche sulle derivate di ordine superiore:
applicandolo, ad esempio, alla derivata parziale $f_{x}$, si trova $f_{xxy}=f_{xyx}$, purché
tali derivate terze siano continue. Pertanto, se una funzione di due variabili ha
derivate parziali continue fino al terzo ordine, le derivate terze distinte sono al più
quattro: $f_{xxx}$, $f_{xxy}$,
$f_{xyy}$, $f_{yyy}$. Conta soltanto quante volte si deriva rispetto a ciascuna variabile,
non l'ordine.

\begin{idea}[L'equazione di Laplace e le correnti potenziali]
In aerodinamica una corrente piana incompressibile e irrotazionale si descrive con un
\emph{potenziale di velocità} $\varphi(x,y)$, le cui derivate parziali sono le componenti
della velocità: $u=\varphi_{x}$, $v=\varphi_{y}$. La conservazione della massa impone
$u_{x}+v_{y}=0$, cioè
\begin{equation}\label{eq:laplace}
  \varphi_{xx}+\varphi_{yy}=0 ,
\end{equation}
l'\emph{equazione di Laplace}; le sue soluzioni si dicono \emph{funzioni armoniche}.
L'irrotazionalità, $v_{x}=u_{y}$, è invece automatica: è proprio il teorema di Schwarz,
$\varphi_{yx}=\varphi_{xy}$. Sono armoniche, per $(x,y)\neq(0,0)$, il potenziale di una
\emph{sorgente} $\log\sqrt{x^{2}+y^{2}}$ (che è, a meno del fattore $\frac12$, la prima delle funzioni
\eqref{eq:radiali}), quello di un \emph{vortice} $\operatorname{arctg}\frac yx$ (nei punti
con $x\neq0$) e quello di una \emph{doppietta} $\frac{x}{x^{2}+y^{2}}$. Sommando la corrente uniforme $\varphi=x$ e la
doppietta si ottiene la corrente attorno a un cilindro circolare: è il primo passo della
teoria classica dei profili alari, che conduce al teorema di Kutta--Žukovskij sulla
portanza. Il programma verifica l'equazione di Laplace con \texttt{sympy}.
\end{idea}

\begin{python}[Funzioni armoniche delle correnti potenziali]
import sympy as sp

x, y = sp.symbols("x y", real=True)
lap = lambda u: sp.simplify(sp.diff(u, x, 2) + sp.diff(u, y, 2))
for nome, phi in (("sorgente", sp.log(sp.sqrt(x**2 + y**2))),
                  ("vortice", sp.atan(y / x)),
                  ("doppietta", x / (x**2 + y**2)),
                  ("cilindro", x + x / (x**2 + y**2)),
                  ("non armonica", x**2 + y**2)):
    print(f"{nome:13} phi = {str(phi):24}  laplaciano = {lap(phi)}")
\end{python}
\begin{risultato}
sorgente      phi = log(sqrt(x**2 + y**2))    laplaciano = 0
vortice       phi = atan(y/x)                 laplaciano = 0
doppietta     phi = x/(x**2 + y**2)           laplaciano = 0
cilindro      phi = x + x/(x**2 + y**2)       laplaciano = 0
non armonica  phi = x**2 + y**2               laplaciano = 4
\end{risultato}

\begin{notastorica}[Eulero, Clairaut, Schwarz, Peano]
L'uguaglianza delle derivate miste fu data per scontata da Nicolaus Bernoulli già nel
1721, e usata da Leonhard Euler in una memoria sulle famiglie di curve scritta nel 1734 (e
pubblicata nel 1740). Alexis Clairaut, nelle ricerche sul calcolo integrale del
1739--1740, se ne servì per un problema che abbiamo appena incontrato: riconoscere quando
un'espressione $M\,dx+N\,dy$ è il differenziale di una funzione $\varphi$. Se lo è, si ha
$M=\varphi_{x}$ e $N=\varphi_{y}$, e l'uguaglianza delle derivate miste impone
$M_{y}=N_{x}$; per una corrente piana, con $M=u$ e $N=v$, è la condizione di
irrotazionalità $u_{y}=v_{x}$. Eulero raccolse questi risultati, insieme al teorema sulle
funzioni omogenee (\cref{subsec:funzioni-omogenee}), nel capitolo sulle funzioni di più
variabili delle \emph{Institutiones calculi differentialis} (1755). Le loro dimostrazioni,
e quelle che seguirono per oltre un secolo, da Lagrange (1797) a Cauchy (1823), erano
incomplete. Nel 1867 Lorenz Leonard Lindelöf,
analizzando sistematicamente le dimostrazioni precedenti, mostrò con un controesempio che
una qualche ipotesi è necessaria; la prima dimostrazione rigorosa, con ipotesi di
continuità del tipo di quelle del \cref{teo:schwarz}, fu data da Hermann Amandus Schwarz
nel 1873. Il controesempio oggi classico, $\frac{xy(x^{2}-y^{2})}{x^{2}+y^{2}}$, dello stesso
tipo di quello dell'\cref{oss:schwarz-controesempio}, fu presentato da Giuseppe Peano nel
1884, nelle note al \emph{Calcolo differenziale} di Genocchi.
\end{notastorica}

\subsection{Approfondimento: un primo sguardo alle equazioni alle derivate parziali}
\label{subsec:edp-primo-sguardo}
\index{equazione!alle derivate parziali}

L'equazione di Laplace \eqref{eq:laplace} è il primo esempio che incontriamo di
\emph{equazione alle derivate parziali}: una relazione tra una funzione incognita di più
variabili e le sue derivate parziali, che si dice di ordine $k$ se vi compaiono derivate
fino all'ordine $k$. Accanto a quella di Laplace, la fisica matematica ha altre due
equazioni modello del secondo ordine, nelle quali una delle variabili è il tempo $t$:
\begin{align}
  u_{t}&=k\,u_{xx} &&\text{(equazione del calore)},\label{eq:calore}\\
  u_{tt}&=c^{2}u_{xx} &&\text{(equazione delle onde)},\label{eq:onde}
\end{align}
con $k>0$ e $c>0$ costanti. La prima descrive la temperatura $u(x,t)$ nel punto di
ascissa $x$ di una sbarra omogenea, all'istante $t$; la costante $k$ è la diffusività
termica del materiale. La seconda descrive lo spostamento $u(x,t)$ di una corda vibrante,
oppure la perturbazione di pressione di un'onda sonora in un condotto; vedremo che $c$ è
la velocità con cui le perturbazioni si propagano. Con gli strumenti di questo capitolo
possiamo verificare che alcune funzioni esplicite le risolvono, e capire da esse
\emph{che cosa} le equazioni descrivono.

\needspace{9\baselineskip}
\begin{proposizione}[Soluzione di d'Alembert dell'equazione delle onde]\label{prop:dalembert}
\index{equazione!delle onde}
Siano $F$ e $G$ funzioni di una variabile derivabili due volte in $\R$, e sia $c>0$.
Allora la funzione
\begin{equation}\label{eq:dalembert}
  u(x,t)=F(x-ct)+G(x+ct)
\end{equation}
soddisfa l'equazione delle onde \eqref{eq:onde} in ogni punto $(x,t)\in\R^{2}$.
\end{proposizione}

\begin{schema}
\Ipotesi{$F$, $G$ derivabili due volte in $\R$; \ $c>0$.}
\Tesi{$u(x,t)=F(x-ct)+G(x+ct)$ verifica $u_{tt}=c^{2}u_{xx}$ in $\R^{2}$.}
\end{schema}

\begin{strategia}
Ogni derivata parziale è la derivata di una funzione di una variabile composta con la
funzione lineare $x-ct$ (oppure $x+ct$): basta la regola della catena del
\cref{teo:derivata-composta}. La funzione interna ha derivata $1$ rispetto a $x$ e $\mp c$
rispetto a $t$; derivando due volte, il fattore $\mp c$ compare al quadrato e produce il
$c^{2}$ dell'equazione.
\end{strategia}

\begin{dimostrazione}
Fissato $t$, la funzione $x\mapsto F(x-ct)$ è la composta di $F$ con la funzione
$x\mapsto x-ct$, che ha derivata $1$; per il \cref{teo:derivata-composta}
$\frac{\partial}{\partial x}F(x-ct)=F'(x-ct)$ e, derivando ancora,
$\frac{\partial^{2}}{\partial x^{2}}F(x-ct)=F''(x-ct)$. Fissato $x$, la funzione
$t\mapsto x-ct$ ha derivata $-c$, quindi
\[
  \frac{\partial}{\partial t}F(x-ct)=-c\,F'(x-ct), \qquad
  \frac{\partial^{2}}{\partial t^{2}}F(x-ct)=c^{2}F''(x-ct).
\]
Allo stesso modo, con $+c$ al posto di $-c$, le derivate seconde di $G(x+ct)$ rispetto a
$x$ e a $t$ valgono $G''(x+ct)$ e $c^{2}G''(x+ct)$. Sommando,
$u_{tt}=c^{2}\bigl[F''(x-ct)+G''(x+ct)\bigr]=c^{2}u_{xx}$.
\end{dimostrazione}

Il significato della \eqref{eq:dalembert} è trasparente: il grafico di $x\mapsto F(x-ct)$
è quello di $F$ traslato di $ct$ verso destra, quindi $F(x-ct)$ è un'onda che viaggia
verso destra con velocità $c$ senza deformarsi, e $G(x+ct)$ un'onda che viaggia verso
sinistra. Ad esempio, $u=\frac12\bigl[F(x-ct)+F(x+ct)\bigr]$ vale $F(x)$ per $t=0$ e ha
$u_{t}(x,0)=\frac12\bigl[-cF'(x)+cF'(x)\bigr]=0$: è la corda che, deformata con il profilo $F$,
viene lasciata andare da ferma. Il profilo si divide in due metà uguali, che si
allontanano in versi opposti (\cref{fig:edp-onde-calore}, a sinistra).

\begin{figure}[htb]
\centering
\includegraphics[width=\textwidth]{edp_onde_calore}
\caption{A sinistra, la soluzione $u=\frac12[F(x-ct)+F(x+ct)]$ dell'equazione delle onde,
con $F(x)=e^{-4x^{2}}$ e $c=1$: il profilo iniziale si divide in due onde che viaggiano
senza deformarsi in versi opposti. A destra, la soluzione \eqref{eq:gaussiana-calore}
dell'equazione del calore, con $k=1$: la campana si abbassa e si allarga
(\texttt{python/fig\_edp\_onde\_calore.py}).}
\label{fig:edp-onde-calore}
\end{figure}

\index{equazione!del calore}Per l'equazione del calore la verifica è altrettanto diretta. La funzione
$u=e^{-kn^{2}t}\sen(nx)$, con $n\in\N$, ha $u_{t}=-kn^{2}u$ e $u_{xx}=-n^{2}u$, e quindi
soddisfa la \eqref{eq:calore}: un profilo sinusoidale di temperatura si smorza
esponenzialmente, tanto più in fretta quanto più è fitto. Più istruttiva è la funzione,
definita per $t>0$,
\begin{equation}\label{eq:gaussiana-calore}
  u(x,t)=\frac{1}{\sqrt{4\pi kt}}\,e^{-x^{2}/(4kt)},
\end{equation}
che il programma più avanti verifica. Per $t$ piccolo è una campana alta e stretta: il
calore depositato in un punto della sbarra all'istante iniziale. Col tempo la campana si
abbassa e si allarga, con una larghezza proporzionale a $\sqrt{kt}$
(\cref{fig:edp-onde-calore}, a destra): là un profilo viaggia senza deformarsi, qui si
smussa e si diffonde.

\begin{esempio}[Le onde di Mach]\label{es:onde-mach}
Un profilo alare sottile sia investito da una corrente supersonica, con numero di Mach
$M>1$. Se le perturbazioni prodotte dal profilo sono piccole, il loro potenziale di
velocità $\varphi(x,y)$, con l'asse $x$ nella direzione della corrente, soddisfa l'equazione
linearizzata
\begin{equation}\label{eq:mach-linearizzata}
  (M^{2}-1)\,\varphi_{xx}-\varphi_{yy}=0 .
\end{equation}
Posto $\beta=\sqrt{M^{2}-1}$, è l'equazione delle onde \eqref{eq:onde} con $y$ al posto di $t$
e $\beta$ al posto di $c$; per la \cref{prop:dalembert} è risolta da
$\varphi=F(x-\beta y)+G(x+\beta y)$. Sopra il profilo, dove le perturbazioni sono portate a
valle, si ha $\varphi=F(x-\beta y)$: la perturbazione è costante lungo le rette
$x-\beta y=\text{costante}$, inclinate sulla direzione della corrente di un angolo $\mu$ con
$\tg\mu=\frac1\beta$, cioè $\sen\mu=\frac1M$. Sono le \emph{onde di Mach}, e $\mu$ è
l'\emph{angolo di Mach}: per $M=2$ si ha $\mu=30^{\circ}$. In regime subsonico, $M<1$,
l'equazione diventa $(1-M^{2})\varphi_{xx}+\varphi_{yy}=0$, e la sostituzione
$\varphi(x,y)=\Phi\bigl(x,\sqrt{1-M^{2}}\,y\bigr)$ la riconduce all'equazione di Laplace
$\Phi_{xx}+\Phi_{yy}=0$ (è la base della regola di Prandtl--Glauert): non ci sono più
direzioni privilegiate, e il profilo «si fa sentire» in tutto il campo, anche a monte.
\end{esempio}

\needspace{7\baselineskip}
\begin{python}[Soluzioni esplicite delle equazioni delle onde e del calore]
import math
import sympy as sp

x, y, t, n = sp.symbols("x y t n", real=True)
k, c, M = sp.symbols("k c M", positive=True)
F, G, d = sp.Function("F"), sp.Function("G"), sp.diff
onde = lambda u: sp.simplify(d(u, t, 2) - c**2 * d(u, x, 2))
calore = lambda u: sp.simplify(d(u, t) - k * d(u, x, 2))
print("u_tt - c^2 u_xx per F(x-ct)+G(x+ct):",
      onde(F(x - c*t) + G(x + c*t)))
print("u_t - k u_xx per e^(-kn^2 t)sen(nx):",
      calore(sp.exp(-k * n**2 * t) * sp.sin(n*x)))
gauss = sp.exp(-x**2 / (4*k*t)) / sp.sqrt(4 * sp.pi * k * t)
print("u_t - k u_xx per la gaussiana:      ", calore(gauss))
# onde di Mach: (M^2 - 1) phi_xx - phi_yy = 0 per phi = F(x - beta y)
beta = sp.sqrt(M**2 - 1)
phi = F(x - beta * y)
print("(M^2-1) phi_xx - phi_yy:",
      sp.simplify((M**2 - 1) * d(phi, x, 2) - d(phi, y, 2)))
for Mach in (1.2, 1.5, 2.0, 3.0):
    mu = math.degrees(math.asin(1 / Mach))
    print(f"M = {Mach:.1f}:  angolo di Mach = {mu:5.2f} gradi")
\end{python}
\begin{risultato}
u_tt - c^2 u_xx per F(x-ct)+G(x+ct): 0
u_t - k u_xx per e^(-kn^2 t)sen(nx): 0
u_t - k u_xx per la gaussiana:       0
(M^2-1) phi_xx - phi_yy: 0
M = 1.2:  angolo di Mach = 56.44 gradi
M = 1.5:  angolo di Mach = 41.81 gradi
M = 2.0:  angolo di Mach = 30.00 gradi
M = 3.0:  angolo di Mach = 19.47 gradi
\end{risultato}

Al crescere di $M$ l'angolo di Mach diminuisce e le onde si ripiegano all'indietro.

\begin{osservazione}[Tre equazioni modello e la fluidodinamica computazionale]
\label{oss:edp-cfd}
Le equazioni di Laplace, del calore e delle onde sono i prototipi di tre famiglie, dette
rispettivamente \emph{ellittiche}, \emph{paraboliche} e \emph{iperboliche}, dai
comportamenti molto diversi. Nelle prime ogni punto risente di tutto il contorno: è il
caso delle correnti subsoniche a potenziale, che i \emph{metodi a pannelli} calcolano
distribuendo sulla superficie del velivolo sorgenti, vortici e doppiette come quelli
dell'idea «L'equazione di Laplace e le correnti potenziali». Le seconde descrivono la
diffusione: nei fluidi è la quantità di moto a diffondere, con la viscosità cinematica
$\nu$ al posto di $k$, e lo spessore della zona rallentata vicino a una parete cresce come
$\sqrt{\nu t}$, come la larghezza della campana \eqref{eq:gaussiana-calore}; è il primo
passo verso la teoria dello strato limite. Le terze descrivono propagazioni a velocità
finita lungo direzioni privilegiate, come l'acustica e le correnti supersoniche
dell'\cref{es:onde-mach}. Le equazioni di Navier--Stokes, risolte dai codici di CFD,
mescolano i tre comportamenti, e un buon metodo numerico rispetta la natura di ciascuno:
nelle regioni iperboliche, ad esempio, gli schemi \emph{upwind} prendono l'informazione
dalla direzione da cui arrivano le onde.
\end{osservazione}

\begin{notastorica}[Le equazioni della fisica matematica]
Le equazioni alle derivate parziali compaiono, ancora sporadiche, nella prima metà del
Settecento. Nella stessa memoria del 1747 in cui scrisse l'equazione della corda vibrante
(si veda la nota storica del \S\ref{sec:derivate-parziali}), d'Alembert ne diede la
soluzione come somma di due onde che viaggiano in versi opposti, cioè la
\cref{prop:dalembert}; il dibattito che ne seguì con Eulero e con Daniel Bernoulli su quali
funzioni $F$ e $G$ fossero ammissibili pesò a lungo sul concetto stesso di funzione
(nota storica del \S\ref{sec:continue}). L'equazione di
Laplace nacque invece dai fluidi: Eulero la scrisse per il potenziale di velocità di una
corrente incomprimibile nei \emph{Principia motus fluidorum}, presentati all'Accademia di
Berlino nel 1752; Pierre-Simon Laplace la ritrovò nel 1782 per il potenziale
gravitazionale, e Siméon-Denis Poisson, nel 1813, mostrò che dentro la materia va corretta
con un secondo membro proporzionale alla densità. Con Fourier, che nel 1807 presentò
l'equazione del calore, con Laplace e con Poisson, la scuola francese fece delle equazioni
alle derivate parziali uno dei campi principali dell'analisi; a Cauchy e a Sof'ja
Kovalevskaja (1874) si deve un teorema di esistenza delle soluzioni. Riemann,
infine, osservò che una funzione armonica con valori assegnati sul bordo di un dominio si
ottiene minimizzando un integrale (il «principio di Dirichlet»); ma che il minimo esista
non è affatto ovvio, e lo dimostrò soltanto David Hilbert, nel 1900.
\end{notastorica}
